\documentclass{homework}
% \usepackage{graphicx} % Required for inserting images
\usepackage{amsmath}
\usepackage{gensymb}
\usepackage{subfig}
\usepackage{float}

\usepackage{algorithm}
\usepackage{algpseudocode}

\class{APPM 5560}
\title{Homework 3}
\author{Sean Jennings}
\date{February 6, 2024}

\begin{document}
\maketitle

\section{Problem 1}
Suppose that there exists some $C>0$ such that $f(x) \leq Cg(x)$ for
all $x\in\R$. Then, for all $x>0$ we have:
\begin{equation*}
    \frac{f(x)}{g(x)} = \frac{\sqrt{2\pi}}{2} e^{-x + x^2/2}
        = \frac{\sqrt{2\pi}}{2} e^{1/2 (x-1)^2 - 1/2} \leq C
\end{equation*}

Solving for $x$ yields:
\begin{align*}
    &e^{1/2(x-1)^2 - 1/2} \leq \frac{2C}{\sqrt{2\pi}} \\
    &\Leftrightarrow 1/2 (x-1)^2 - 1/2 \leq \ln(\frac{2C}{\sqrt{2\pi}}) \\
    &\Leftrightarrow x \leq 1 + \sqrt{1 + 2 \ln(\frac{2C}{\sqrt{2\pi}})}
\end{align*}
This contradicts that $x$ can be any real number and thus no such
number $C$ can exist.

\section{Problem 2}
\begin{letterenum}
\item The CDF of $X$, for $-\pi/2 \leq x \leq \pi/2$ is given by:
\begin{equation*}
    F_X(x) = \int_{-\pi/2}^{x} \cos(y)/2 dy 
        = 1/2 (\sin(y))\eval_{-\pi/2}^x = \frac{1}{2} (1 + \sin(x))
\end{equation*}

Thus, we can compute $X$ via the inverse transform method by the following
relations:
\begin{align*}
    U = F(X) &= 1/2(1 + \sin(x)) \\
    \Leftrightarrow X &= \sin^{-1}(2U - 1)
\end{align*}

In pseudocode:
\begin{algorithmic}
    \State $U \sim$  Uniform(0, 1)
    \State \Return $\sin^{-1}(2U - 1)$
\end{algorithmic}

\item Since the PDF of $X$ can be bounded by a box of width $\pi$ and height
$1/2$, the acceptance-rejection method can be performed by drawing two uniform
random variables, $U_1$ and $U_2$, and accepting if the point $(U_1, U_2)$ is
under the PDF curve: $(X, f(X))$.

\begin{algorithmic}
    \Loop
        \State $U_1 \sim$ Uniform(0, 1)
        \State $U_2 \sim$ Uniform(0, 1)
        \State $X \gets \pi \cdot (U_1 - 1/2)$
        \If {$U_2 < \cos(X)$}
            \Comment Equivalent to $Y=1/2 U_2$ and $Y < \cos(X)/2$
            \State \Return $X$
        \EndIf
    \EndLoop
\end{algorithmic}
\end{letterenum}

\section{Problem 3}
\begin{letterenum}
\item We define $h: \R \to \R$ to be the function:
\begin{align*}
    h(x) &= \frac{f(x)}{g(x)} = \frac{e^{-x^2/2}}{\sqrt{2\pi}} \bigl(\pi(1 + x^2)\bigr) \\
        &= \sqrt{\frac{\pi}{2}} (1 + x^2) e^{-x^2/2}
\end{align*}

Taking the derivative yields:
\begin{align*}
    \frac{dh}{dx}(x) = \sqrt{\pi/2} \biggl(
        2x e^{-x^2/2} - x(1 + x^2) e^{-x^2/2}
    \biggr)
    &= \sqrt{\pi/2} x e^{-x^2/2} (1 - x^2)
\end{align*}
Equating $dh/dx$ to zero, we find that $h(x)$ has extrema at $x=0,\pm 1$.
Since $e^{-x^2/2}$ decreases faster than $(1+x^2)$ increases, $h(x)$ approaches
0 for large $x$. Furthermore,
\begin{equation*}
    h(0) = \sqrt{\pi/2} < \sqrt{2\pi}e^{-1/2} = h(\pm1).
\end{equation*}
We take $C$ to be the maximum value of $h$:
\begin{equation*}
    C = \max_{x\in\R} h(x) = h(1) = \sqrt{2\pi}e^{-1/2}
\end{equation*}

\item The CDF of $Y$ is given by
\begin{align*}
    G_Y(y) &= \int_{-\infty}^{y} \frac{1}{\pi(1 + x^2)} dx \\
        &= \frac{1}{\pi}(\tan^{-1}x) \eval_{x=-\infty}^y \\
        &= \frac{1}{\pi} \biggl(\tan^{-1}y + \pi/2 \biggr)
\end{align*}
Inverting the equation $U = G(Y)$ we have:
\begin{align*}
    &U = G(Y) = 1/2 + 1/\pi \tan^{-1} y \\
    &\Leftrightarrow y = \tan \biggl(\pi \cdot (u - 1/2)\biggr)
\end{align*}
And in pseudocode:
\begin{algorithmic}
    \Function{Lorentzian}{}
        \State $U \sim $ Uniform(0, 1)
        \State \Return $\tan (\pi \cdot(U - 1/2))$
    \EndFunction
\end{algorithmic}

\item To simulate the Gaussian, sample from the Lorentzian and accept
if $U \sim $ Uniform(0,1) is less than:
\begin{equation*}
    \frac{f(x)}{Cg(x)} = h(x) / C = \frac{\sqrt{\pi/2} e^{-x^2/2}( 1+ x^2)}
                                          {(\sqrt{2\pi} e^{-1/2})}
        = \frac{1}{2} e^{1/2(1 - x^2)}(1 + x^2)
\end{equation*}

In pseudocode:
\begin{algorithmic}
    \Loop
        \State $U \sim$ Uniform(0, 1)
        \State $X \sim $ Lorentzian()
        \If {$U < 1/2 \cdot (1 + x^2) \cdot e^{1/2(1 - x^2)}$}
            \State \Return $X$
        \EndIf
    \EndLoop
\end{algorithmic}
\end{letterenum}

\section{Problem 4}
Consider these two events
\begin{enumerate}
    \item $A = \{X_0=0, X_1=1\}$ Red on draw 0, white of draw 1,
    \item $B = \{X_0=1, X_1=1\}$ White on draw 0, white on draw 1
\end{enumerate}

Since the balls are drawn without replacement, they depend not only on the color
of the last ball drawn, but on the the number
of balls remaining of each color in the jar. For event $A$, 49 out of the 98
remaining balls are white so:
\begin{equation*}
    {\bf P}(X_2 = 1 | A) = 49/98
\end{equation*}
while for $B$, only 48 out of 98 are white and we have
\begin{equation*}
    {\bf P}(X_2 = 1 | B) = 48/98
\end{equation*}
Clearly, we see that $X_t$ does not satisfy the Markovian property
since $X_2$ depends $X_0$. I.e.

\begin{equation*}
    {\bf P}(X_2=1 | X_1 = 1 \cap X_0 = 1) \neq {\bf P}(X_2=1 | X_1 = 1 \cap X_0 = 0).
\end{equation*}

\section{Problem 5}
\begin{letterenum}
\item Since $X_t$, the number of working lightbulbs, is an integer-valued variable
with $0 \leq X_t \leq 100$, the sample space $S$ is the integers from 0 to 100.

\item This is a Markov chain, since the number of lightbulbs remaining in month
$t+1$ depends only upon the number of lightbulbs remaining in the previous month.
Specifically,
the number of lightbulbs remaining at month $t+1$ is a Binomial with $X_t$ trials of probability $1-p$ (the probability each lightbulb remains is 1 minus the probability
that it fails).

\item The transition probabilities given by the distribution Binomial($X_t$, $1-p$):
\begin{equation*}
    {\bf P}(X_{t+1} = i | X_t = j) = \choose{j}{i} (1 - p)^i p^{j-i}
\end{equation*}

\section{Problem 6}
The possible paths that lead to $X_4=1$ are:
\begin{itemize}
    \item A: $3\to2\to1\to1\to1$
    \item B: $3\to2\to3\to2\to1$
    \item C: $3\to4\to3\to2\to1$
\end{itemize}

The total probability of a path is the product of each step, so
the probability of each individual path is given by:

\begin{align*}
    {\bf P}(A) &= p(3,2)p(2,1)p(1,1)p(1,1) \\
        &= 1/4 \cdot 1/4 \cdot 1 \cdot 1 = 1/16 \\
    {\bf P}(B) &= p(3,2)p(2,3)p(3,2)p(2,1) \\
        &= 1/4 \cdot 3/4 \cdot 1/4 \cdot 1/4 = 3/256 \\
    {\bf P}(C) &= p(3,4)p(4,3)p(3,2)p(2,1) \\
        &= 3/4 \cdot 1/4 \cdot 1/4 \cdot 1/4 = 3/256
\end{align*}

and thus the total probability of going from 3$\to$1 in 4 steps is the sum
of these (disjoint) events:
\begin{equation*}
    {\bf P}(3 \to 1 \text{ in 4 steps}) = {\bf P}(A) + {\bf P}(B) + {\bf P}(C)
        = 8/128 + 2(3/256) = 11/128
\end{equation*}

\item The state transition matrix is given by:
\begin{equation*}
    A = \mat{
        1 & 0 & 0 & 0 & 0 \\
        1/4 & 0 & 3/4 & 0 & 0 \\
        0 & 1/4 & 0 & 3/4 & 0 \\
        0 & 0 & 1/4 & 0 & 3/4 \\
        0 & 0 & 0 & 0 & 1
    }
\end{equation*}

\item Taking the matrix power (using Python) we have:
\begin{equation*}
    A^4 = \frac{1}{128} \mat{
        128 & 0 & 0 & 0 & 0 \\
        38 & 9 & 0 & 27 & 54 \\
        11 & 0 & 18 & 0 & 99 \\
        2 & 3 & 0 & 9 & 114 \\
        0 & 0 & 0 & 0 & 128
    }
\end{equation*}
The probability computed by summing individual paths from $X_0=3$ to $X_4=1$ indeed
corresponds to the (3,1)-th entry of $A^4$.

\item $A^100$ is given by:
\item \begin{equation*}
    A^100 = \mat{
        1 & 0 & 0 & 0 & 0 \\
        0.325 & 0 & 0 & 0 & 0.675 \\
        0.1 & 0 & 0 & 0 & 0.9 \\
        0.025 & 0 & 0 & 0 & 0.975 \\
        0 & 0 & 0 & 0 & 1
    }
\end{equation*}
so probability of eventually ending in state 1 from state 3 is again given by
the (3,1)-th entry which is $A^100(3,1) = 0.1$.
\end{letterenum}

\section{Problem 7}
\begin{letterenum}
\item $X_t$, the sum of $x$ and $y$ coordinates, takes its minimum of 0
at the point $(0, 0)$ and its maximum of 4 at point $(2,2)$, and the
other points on the grid correspond to the integers between. Thus, the sample
space $S$ is the set $\{0, 1, \dots, 4\}$.

\item The transition probability matrix is:
\begin{equation*}
    A = \mat{
        0 & 1 & 0 & 0 & 0 \\
        1/3 & 0 & 2/3 & 0 & 0 \\
        0 & 1/2 & 0 & 1/2 & 0 \\
        0 & 0 & 2/3 & 0 & 1/3 \\
        0 & 1 & 0 & 1 & 0
    }
\end{equation*}

\item There is only one unique shortest path: $0\to1\to2\to3\to4$ and its
probability is: 
\begin{equation*}
    {\bf P}(0\to1\to2\to3\to4) = p(0, 1)p(1,2) p(2,3) p(3,4)
        = 1 \cdot 2/3 \cdot 1/2 \cdot 1/3 = 1/9
\end{equation*}

\end{letterenum}

\section{Problem 8}

\begin{letterenum}
\item The transition matrix is:
\begin{equation*}
    A = \mat{
        0 & 0.5 & 0.5 \\
        0.1 & 0.7 & 0.2 \\
        0.2 & 0.8 & 0
    }
\end{equation*}

\item Next sunday is 7 steps within the markov chain, so we compute
the multi-step transition probability by taking the transition matrix to the 7th
power and looking at the (1,1)-th entry:
\begin{equation*}
    A^7(1,1) \approx 0.1084
\end{equation*}

\item The transition matrix to the 50th power is:
\begin{equation*}
    A^50 \approx \mat{
        0.1085 & 0.6977 & 0.1938 \\
        \vdots & \vdots & \vdots
    }
\end{equation*}
where the dots indicate that each row is identical (i.e. the initial state is
irrelevant). Thus, as $t\to\infty$ the probability that dish 2 will be eaten
on day $t$ is about 0.6977.
\end{letterenum}





\end{document}
